
% \begin{center}
% Министерство цифрового развития, связи и массовых коммуникаций Российской Федерации\\
% федеральное государственное бюджетное образовательное учреждение высшего образования\\
% «Сибирский государственный университет телекоммуникаций и информатики»\\
% (СибГУТИ)
% \end{center}

% \vspace{0.6cm}

% \begin{flushright}
% \underline{09.03.01 Информатика и вычислительная техника}\\
% {\footnotesize (направление подготовки/специальность)}\\[0.3em]
% \underline{Программное обеспечение мобильных систем}\\
% {\footnotesize (профиль/специализация)}\\[0.3em]
% \underline{Очная}\\
% {\footnotesize (форма обучения)}
% \end{flushright}

% \vspace{1cm}

% \begin{center}
% \textbf{ОТЧЕТ ПО ПРОИЗВОДСТВЕННОЙ ПРАКТИКЕ}\\[0.2em]
% {\footnotesize (вид практики)}\\[0.7em]

% \textbf{Тип практики}\\[0.4em]
% \underline{Технологическая (проектно-технологическая) практика на предприятии ООО «Бюро 1440»}\\
% {\footnotesize (наименование профильной организации/ структурного подразделения СибГУТИ)}\\[1em]

% \textbf{ТЕМА ИНДИВИДУАЛЬНОГО ЗАДАНИЯ}\\[0.3em]
% \underline{TODO}
% \end{center}

% \vspace{0.5cm}

% \begin{flushright}
% \textbf{Выполнил:}\\[0.2em]
% студент института информатики и вычислительной техники\\
% Гмыря Ярослав Александрович\\
% группа \underline{ИА-331}
% \end{flushright}

% \vspace{1.2cm}

% «\underline{\hspace{1cm}}» \underline{\hspace{2.5cm}} 202\underline{\hspace{0.7cm}} г. \hfill
% \underline{\hspace{4cm}} / Гмыря Я.А. /

% \vspace{0.25cm}
% \hfill {\footnotesize (подпись)\hspace{1.2cm}(ФИО)}

% \vspace{1.1cm}

% \textbf{Проверил\textsuperscript{1}:}\\
% Руководитель практики от профильной организации

% \vspace{0.5cm}

% «\underline{\hspace{1cm}}» \underline{\hspace{2.5cm}} 202\underline{\hspace{0.7cm}} г. \hfill
% \underline{\hspace{4cm}} / Андреев А.В. /

% \vspace{0.25cm}
% \hfill {\footnotesize (подпись)\hspace{1.2cm}(ФИО)}

% \vspace{1.1cm}

% \textbf{Проверил:}\\
% Руководитель практики от СибГУТИ

% \vspace{0.5cm}

% «\underline{\hspace{1cm}}» \underline{\hspace{2.5cm}} 202\underline{\hspace{0.7cm}} г. \hfill
% \underline{\hspace{4cm}} / Брагин К.И. /

% \vspace{0.25cm}
% \hfill {\footnotesize (подпись)\hspace{1.2cm}(ФИО)}

% \vspace{1.1cm}

% отметка\textsuperscript{2} \underline{\hspace{6cm}}\\[0.7em]
% \hfill «\underline{\hspace{1cm}}» \underline{\hspace{2.5cm}} 202\underline{\hspace{0.7cm}} г.

% \vspace{0.8cm}

% \begin{center}
% Новосибирск 2025
% \end{center}

% \vspace{0.4cm}

% \noindent\rule{6cm}{0.4pt}

% {\footnotesize
% \textsuperscript{1} В случае прохождения практики в профильной организации\\
% \textsuperscript{2} Заполняется во время промежуточной аттестации
% }

% \vfill
% \newpage

% \begin{center}
% \textbf{План-график проведения}\\[0.2em]
% \textbf{Производственной практики}\\[0.3em]
% {\footnotesize (вид практики)}
% \end{center}

% \vspace{0.5cm}

% \underline{Гмыря Ярослав Александрович}\\
% {\footnotesize Фамилия Имя Отчество студента}

% \vspace{0.3cm}

% Институт ИВТ, курс 3, гр. \underline{ИА-331}\\[0.2em]
% Направление: \underline{09.03.01 Информатика и вычислительная техника}\\
% {\footnotesize Код – Наименование направления (специальности)}

% \vspace{0.4cm}

% Направленность (профиль): \underline{Программное обеспечение мобильных систем}\\
% Место прохождения: \underline{ООО «Бюро 1440», Новосибирск}\\
% Адрес практики: 314\\
% Объем практики: 360 / 10\\
% Срок практики: \textbf{16.09.2025}

% \vspace{0.6cm}

% Содержание практики:\\
% Тема индивидуального задания практики: \underline{TODO}

% \vspace{0.6cm}

% \begin{center}
% \begin{tabular}{|p{10.8cm}|c|}
% \hline
% \textbf{Наименование видов деятельности практики} & \textbf{Дата (начало–окончание)}\\

% \hline
% Архитектура Adalm Pluto SDR. GNU Radio. Построение радио-приёмника & 16.09.2025\\

% \hline
% Введение в архитектуру SDR-устройств. Знакомство с библиотеками Soapy SDR, Libio для работы с Adalm Pluto SDR. 
% Инициализация SDR-устройства. Работа с буфером: получение цифровых IQ-отсчетов. & 23.09.2025\\

% \hline
% Принципы работы библиотеки Soapy SDR и работы с Adalm Pluto. Работа с библиотеками Soapy SDR, Libio. 
% Формирование и передача с SDR сигналов произвольной формы & 30.09.2025\\

% \hline
% Архитектура SDR-устройств. Примеры формирования I/Q-сэмплов произвольной формы. Работа с буфером приема SDR & 7.10.2025\\

% \hline
% Прием сигналов с фазовой модуляцией BPSK/QPSK. Имитация аналоговой передачи звука и его прием с использованием SDR. 
% Анализ влияния чувствительности приемника и усиления передатчика на качество принятых отсчетов сигнала (семплов) & 14.10.2025\\

% \hline
% Имитация аналоговой передачи звука и его прием с использованием SDR. Анализ влияния чувствительности приемника и усиления 
% передатчика на качество принятых отсчетов сигнала (семплов) & 21.10.2025\\
% \hline

% \hline
% Моделирование формирования и приема QPSK-сигналов. Реализация приема и передачи BPSK-сигналов & 28.10.2025\\
% \hline

% \hline
% Алгоритм дискретной свертки. Дискретная свертка. Реализация приема и передачи BPSK-символов & 11.11.2025\\
% \hline

% \hline
% Прием QPSK BPSK, прием на согласованный фильтр, глазковая диаграмма, 
% поиск оптимального отсчетного значения и необходимость символьной синхронизации. Прием и фильтрация сигнала. 
% Прямоугольный и приподнятый косинус & 18.11.2025\\
% \hline

% \hline
% Символьная синхронизация, детектор временной ошибки, схема Гарднера. Программная реализация детектора временной ошибки 
% (синхронизация приемника и передатчика) на SDR & 25.11.2025\\
% \hline
% \end{tabular}
% \end{center}

% \vspace{0.8cm}

% В соответствии с рабочей программой практики\\
% Руководитель практики от профильной организации\\
% «\underline{\hspace{1cm}}» \underline{\hspace{2.5cm}} 2025 г. \hfill
% \underline{\hspace{4cm}} / Андреев А.В. /

% \vspace{0.2cm}
% \hfill {\footnotesize (подпись)\hspace{1.2cm}(ФИО)}

% \vspace{0.8cm}

% Руководитель практики от СибГУТИ\\
% «\underline{\hspace{1cm}}» \underline{\hspace{2.5cm}} 2025 г. \hfill
% \underline{\hspace{4cm}} / Брагин К.И. /

% \vspace{0.2cm}
% \hfill {\footnotesize (подпись)\hspace{1.2cm}(ФИО)}

% \vfill

% {\footnotesize
% * В случае прохождения практики в профильной организации.
% }

% \vfill
% \newpage

% % \thispagestyle{empty}
% % \newcommand{\reviewline}{\noindent\rule{\textwidth}{0.3pt}\par\vspace{0.28cm}}

% % \begin{center}
% % \textbf{Отзыв о работе студента}
% % \end{center}

% % \vspace{0.6cm}
% % \reviewline
% % \begin{center}
% % {\footnotesize (ФИО студента)}
% % \end{center}
% % \vspace{0.2cm}

% % \reviewline
% % \reviewline
% % \reviewline
% % \reviewline
% % \reviewline
% % \reviewline
% % \reviewline
% % \reviewline
% % \reviewline
% % \reviewline

% % \vspace{0.5cm}

% % \begin{center}
% % \textbf{Уровень освоения компетенций}
% % \end{center}

% % \vspace{0.4cm}

% % \begin{center}
% % \renewcommand{\arraystretch}{1.2}
% % \begin{tabular}{|p{9cm}|p{5cm}|}
% % \hline
% % \multicolumn{1}{|c|}{\textbf{Компетенции}} &
% % \multicolumn{1}{c|}{\textbf{Уровень сформированности компетенций}}\\
% % \hline
% % \textbf{ПК-1} Способен разрабатывать требования и
% % проектировать программное обеспечение
% % &
% % \\
% % \hline
% % \textbf{ПК-3} Способен осуществлять эксплуатацию и развитие
% % транспортных сетей и сетей передачи данных, включая
% % спутниковые системы
% % &
% % \\
% % \hline
% % \end{tabular}
% % \end{center}

% % \vspace{0.35cm}

% % {\footnotesize Уровень компетенций: высокий, средний, низкий, не сформирована}

% % \vspace{0.7cm}


% % \textbf{Руководитель практики от СибГУТИ:}\\[0.25cm]

% % \noindent
% % \begin{minipage}{0.45\textwidth}
% % Старший преподаватель\\
% % Кафедры ТС и ВС\\[0.2em]
% % {\footnotesize (должность руководителя практики)}
% % \end{minipage}
% % \hfill
% % \begin{minipage}{0.22\textwidth}
% % \centering\underline{\hspace{2.5cm}}\\[0.2em]
% % {\footnotesize (подпись)}
% % \end{minipage}
% % \hfill
% % \begin{minipage}{0.22\textwidth}
% % \centering\underline{Брагин К.И.}\\[0.2em]
% % {\footnotesize (ФИО руководителя)}
% % \end{minipage}

% % \vspace{2.5cm}

% % «\underline{\hspace{1cm}}» \underline{\hspace{2.8cm}} 2025 г.

% % \vfill
% % \newpage

% \begin{center}
%     {\Large \textcolor{blue}{\textbf{Оглавление}}}
% \end{center}

% \vspace{1em}

% \noindent
% Занятие 1 \dotfill 5\\[0.5em]
% Занятие 2 \dotfill 27\\[0.5em]
% Занятие 3 \dotfill 40\\[0.5em]
% Занятие 4 \dotfill 52\\[0.5em]
% Занятие 5 \dotfill 72\\[0.5em]
% Занятие 6 \dotfill 85\\[0.5em]
% Занятие 7 \dotfill 94\\[0.5em]
% Занятие 8 \dotfill 107\\[0.5em]
% Занятие 9 \dotfill 115\\[0.5em]
% Занятие 10 \dotfill 131\\[0.5em]


% \endinput

